% =========================================================
%  sesi/08/sesi.tex
%  Sesi 8 — Grid dan Pencarian Jalur A*
% =========================================================

\sesi{08}{Grid dan Pencarian Jalur A*}{2.3}{150 menit}

Pada minggu 8 penjaga akhirnya memperoleh pengetahuan tentang ruang. Sejak
sesi 7 ia sudah memiliki mesin keadaan: berdiam, berpatroli, dan mengejar.
Namun ketika mengejar, ia hanya berjalan lurus ke arah pemain. Selama jalan
itu kosong, perilakunya tampak cerdas. Begitu sebuah pilar atau dinding
berdiri di antaranya, ia menekan dinding sampai pemain menghilang dari
pandangan. Ia tahu ke mana harus pergi, tetapi tidak tahu bagaimana sampai.

Sesi ini menutup kekurangan itu. Peta Lorong yang selama ini hanya menjadi
gambar di layar dibaca sebagai graf, lalu dicari rutenya dengan algoritma
A*. Pekerjaan mencari rute pada peta berhalangan itu disebut
\istilah{pathfinding}. Penjaga tidak lagi menempuh garis lurus; ia mengikuti
rangkaian petak yang dijamin tidak melewati dinding. Rute itu digambar di
layar agar Anda dapat memeriksa hasil pencarian, bukan hanya mempercayainya.

Kuis 2 pada minggu ini mencakup Materi 6--7, yaitu perilaku mengejar dan
mesin keadaan. Gunakan kuis itu sebagai pemeriksaan sebelum kedua lapisan
disatukan: bila keadaan penjaga belum Anda pahami, pencarian jalur hanya
akan menutupi kekeliruan yang lebih dasar.

% ---------------------------------------------------------
\bagian{Tujuan Sesi}

\begin{tujuanbox}
Setelah menuntaskan sesi ini, Anda dapat menunjukkan:
\begin{enumerate}[leftmargin=1.5em, itemsep=0.2em]
  \item penjaga yang mencapai pemain di balik pilar dengan menempuh jalan
        memutar, bukan menekan dinding;
  \item peta Lorong $24 \times 14$ petak yang diubah menjadi graf: sel
        36 piksel, dinding dibuang, dan setiap sisi berbiaya satu;
  \item perhitungan $g$, $h$, dan $f$ untuk sebuah petak dengan angka yang
        benar, termasuk $h = |\Delta x| + |\Delta y|$;
  \item penjelasan mengapa \code{heapq} mengeluarkan petak dengan $f$
        terkecil dan mengapa isi antrean yang usang harus dibuang;
  \item jalur yang direkonstruksi dari tabel asal dan digambar sebagai garis
        beserta titik pusat setiap petaknya.
\end{enumerate}
\end{tujuanbox}

% ---------------------------------------------------------
\bagian{Prasyarat}

\begin{prasyaratbox}
\begin{itemize}[leftmargin=1.4em, itemsep=0.25em]
  \item Checkpoint sesi 7 berjalan: penjaga berpindah antara keadaan
        \code{idle}, \code{patrol}, dan \code{chase}.
  \item Anda dapat memakai \code{dict}, \code{set}, dan \code{tuple} sebagai
        pasangan koordinat, serta menulis fungsi yang mengembalikan daftar.
  \item Anda memahami pemakaian \code{delta_seconds} dan \code{pygame.Vector2}
        dari sesi 3 dan 6.
  \item Bab X buku ajar sudah dibaca, khususnya bagian tentang graf,
        \istilah{heuristik}, dan rekonstruksi jalur.
\end{itemize}
\end{prasyaratbox}

% ---------------------------------------------------------
\bagian{Peta Kemampuan}

\petasesi
  {Penjaga bermesin keadaan; saat mengejar ia bergerak lurus dan berhenti
   ketika dinding menghalangi.}
  {Peta berpetak sebagai graf, algoritma A* dengan
   \istilah{antrean prioritas}, rekonstruksi jalur, dan jalur yang digambar
   pada layar.}
  {Penggabungan FSM dan A* dengan \istilah{repath} berkala, dua level
   berurutan, dan tiga \istilah{skenario uji} pada sesi 9.}

% ---------------------------------------------------------
\bagian{Konsep yang Dipakai}

\subsection*{Dari peta huruf ke graf}

Peta Lorong adalah daftar 14 baris teks, masing-masing 24 karakter. Setiap
karakter menempati satu petak berukuran 36 piksel. Karakter \texttt{\#}
menyatakan dinding, sedangkan karakter lain menyatakan lantai. Dua petak
bersebelahan secara mendatar atau tegak dihubungkan oleh satu sisi, dan
setiap sisi berbiaya satu karena memindahkan penjaga satu petak selalu
menempuh jarak yang sama.

Perubahan itu menghasilkan tiga hal: himpunan dinding, lebar peta, dan
tinggi peta. Petak dinding tidak pernah menjadi simpul, sehingga pencarian
tidak perlu memeriksa apakah sebuah langkah menembus dinding. Cukup dengan
membuangnya sekali ketika peta dibaca.

\subsection*{Biaya, perkiraan, dan jumlahnya}

Tiga lambang dipakai sepanjang sesi ini.

\begin{description}[leftmargin=1.4em, itemsep=0.15em]
  \item[$g$] biaya dari titik awal sampai petak yang sedang diperiksa. Pada
        peta ini $g$ sama dengan banyaknya langkah yang sudah ditempuh.
  \item[$h$] perkiraan biaya dari petak itu sampai tujuan. Pada gerak empat
        arah dipakai jarak Manhattan,
        \[
          h = |\Delta x| + |\Delta y|,
          \qquad
          \Delta x = x_{\text{tujuan}} - x_{\text{petak}},
          \quad
          \Delta y = y_{\text{tujuan}} - y_{\text{petak}} .
        \]
        Di sini $\Delta x$ adalah selisih kolom dan $\Delta y$ selisih baris.
  \item[$f$] perkiraan biaya total jalur yang melalui petak itu,
        \[
          f = g + h .
        \]
\end{description}

Petak dengan $f$ terkecil diperiksa lebih dahulu. Karena $h$ menghitung sisa
jarak, petak yang menuju tujuan mendapat $f$ lebih kecil daripada petak yang
menjauh. Pencarian menjadi terarah, bukan menyebar ke segala arah seperti
pencarian melebar. Jarak Manhattan tidak pernah melebihkan jarak sebenarnya
pada gerak empat arah; sifat itu yang menjaga jalur hasil pencarian tetap
terpendek.

\subsection*{Antrean prioritas}

Agar petak dengan $f$ terkecil selalu keluar lebih dahulu, tepi pencarian
disimpan sebagai \istilah{antrean prioritas}. Python menyediakannya pada
modul \code{heapq}. Isi antrean berupa rangkap tiga
\code{(f, g, petak)}. Bila dua petak memiliki $f$ sama, angka $g$ menjadi
pembanding kedua, lalu pasangan koordinat petak menjadi pembanding terakhir.
Susunan itu membuat urutan pemeriksaan dapat diulang dengan hasil yang sama.

Sebuah petak boleh masuk antrean lebih dari sekali, yaitu setiap kali
ditemukan jalan yang lebih murah menuju petak itu. Karena itu antrean dapat
memuat isi usang. Program membuangnya dengan dua cara: himpunan
\code{closed} menolak petak yang sudah diperiksa, dan kamus \code{best_g}
menolak jalan yang tidak lebih murah daripada yang sudah diketahui.

\subsection*{Menyimpan asal dan menyusun jalur}

Selama pencarian, \code{came_from} menyimpan dari petak mana setiap petak
ditemukan. Ketika tujuan keluar dari antrean, jalur diperoleh dengan
menelusuri kamus itu mundur sampai titik awal, lalu membalik urutannya.
Apabila tujuan tidak pernah keluar dari antrean, fungsi mengembalikan daftar
kosong. Keadaan itu wajar terjadi, misalnya ketika pintu keluar terkurung
dinding, dan harus ditangani oleh pemanggil.

\subsection*{Menggambar jalur}

Jalur berupa daftar petak, sedangkan layar memakai piksel. Setiap petak
diubah menjadi titik pusatnya, lalu titik-titik itu dihubungkan dengan
\code{pygame.draw.lines}. Bulatan kecil pada setiap titik memperlihatkan
petak mana yang dilalui. Jalur hanya digambar ketika penjaga mengejar,
sehingga garis kuning di layar selalu berarti rute yang sedang dipakai.

% ---------------------------------------------------------
\bagian{Tutorial}

Kerjakan langkah berikut berurutan. Setiap langkah adalah satu tindakan yang
selesai; jalankan program setiap kali satu bagian baru selesai ditulis.

\begin{langkah}
  \item \textbf{Siapkan folder kerja.} Salin folder proyek sesi 7 menjadi
        folder sesi 8, lalu jalankan \code{python main.py --level 2} untuk
        memastikan penjaga bermesin keadaan masih bekerja. Checkpoint lama
        tidak boleh diubah.

  \item \textbf{Tambahkan pengaturan petak.} Pada \file{settings.py},
        pastikan \code{CELL = 36}, \code{GRID_OFFSET_X = 48}, dan
        \code{GRID_OFFSET_Y = 76} tersedia. Tambahkan pengaturan penjaga
        yang dipakai pencarian: \code{GUARD_RADIUS_CELLS = 7.0},
        \code{GUARD_PATROL_SPEED = 90.0}, \code{GUARD_CHASE_SPEED = 165.0},
        dan \code{WAYPOINT_REACHED = 6.0}.

  \item \textbf{Buat berkas pencarian.} Buat
        \file{gudang_artefak/pathfinding.py}. Tulis fungsi
        \code{build_grid} yang menyisir baris peta, mengumpulkan petak
        dinding ke dalam himpunan, lalu mengembalikan dinding, lebar, dan
        tinggi peta.

  \item \textbf{Tulis jarak Manhattan.} Tambahkan fungsi \code{manhattan}
        yang menerima dua pasangan petak dan mengembalikan jumlah selisih
        kolom dan selisih baris. Fungsi ini kecil, tetapi menjadi jantung
        heuristik pada seluruh sesi.

  \item \textbf{Tulis penentuan tetangga.} Tambahkan fungsi
        \code{neighbours} yang memeriksa empat arah, membuang petak di luar
        batas peta, dan membuang petak dinding. Keluaran fungsi ini berupa
        generator, sehingga pemanggil dapat menelusurinya dengan
        \code{for}.

  \item \textbf{Susun kerangka pencarian.} Tulis \code{find_path} dengan
        antrean \code{open_heap} berisi \code{(manhattan(start, goal), 0,
        start)}, kamus \code{came_from}, kamus \code{best_g}, dan himpunan
        \code{closed}. Periksa lebih dahulu apakah titik awal atau tujuan
        berada di dinding; bila ya, kembalikan daftar kosong.

  \item \textbf{Tulis perulangan A*.} Keluarkan isi antrean dengan
        \code{heapq.heappop}. Lewati petak yang sudah ada di \code{closed}.
        Bila petak itu sama dengan tujuan, rekonstruksi jalurnya. Untuk
        setiap tetangga, hitung \code{new_g = g + 1}; bila lebih kecil
        daripada catatan sebelumnya, perbarui \code{best_g},
        \code{came_from}, lalu masukkan \code{(new_g + h, new_g, tetangga)}
        ke antrean.

  \item \textbf{Tulis rekonstruksi jalur.} Tambahkan \code{_reconstruct}
        yang menelusuri \code{came_from} dari tujuan ke titik awal,
        mengumpulkan petaknya, lalu membalik urutan daftar sebelum
        dikembalikan.

  \item \textbf{Sambungkan ke penjaga.} Pada \file{game.py}, ubah
        \code{_chase} agar memanggil \code{find_path} dari petak penjaga ke
        petak pemain, menyimpan hasilnya pada \code{self.path} dan
        \code{self.guard.path}, lalu memanggil \code{self.guard.follow}.
        Tambahkan \code{_player_cell} untuk mengubah titik pusat pemain
        menjadi pasangan petak.

  \item \textbf{Gambar jalurnya.} Pada \code{_draw_lorong}, ubah setiap
        petak jalur menjadi titik pusat piksel, hubungkan dengan
        \code{pygame.draw.lines}, lalu tandai titik pusat itu dengan
        bulatan kecil. Jalankan \code{python main.py --level 2}, lalu
        bersembunyilah di balik pilar dan amati rute penjaga.
\end{langkah}

\begin{center}
% figures/grid-astar.tikz
% Grid dengan dinding, titik awal, tujuan, dan jalur hasil A*.
\begin{tikzpicture}[x=0.62cm, y=0.62cm]
  \foreach \x in {0,...,9}{
    \foreach \y in {0,...,5}{
      \draw[ubLine!60, very thin] (\x,\y) rectangle ++(1,1);
    }
  }
  % dinding
  \foreach \x/\y in {3/1, 3/2, 3/3, 6/2, 6/3, 6/4}{
    \fill[ubNavy!85] (\x,\y) rectangle ++(1,1);
  }
  % jalur
  \foreach \x/\y in {0/0, 1/0, 2/0, 2/1, 2/2, 2/3, 2/4, 3/4, 4/4, 5/4, 5/3, 5/2, 5/1, 5/0, 6/0, 7/0, 8/0, 8/1, 8/2, 9/2}{
    \fill[ubAccent!18] (\x,\y) rectangle ++(1,1);
  }
  \draw[ubAccent, thick, -{Stealth[length=2mm]}]
    (0.5,0.5) -- (1.5,0.5) -- (2.5,0.5) -- (2.5,1.5) -- (2.5,2.5) --
    (2.5,3.5) -- (2.5,4.5) -- (3.5,4.5) -- (4.5,4.5) -- (5.5,4.5) --
    (5.5,3.5) -- (5.5,2.5) -- (5.5,1.5) -- (5.5,0.5) -- (6.5,0.5) --
    (7.5,0.5) -- (8.5,0.5) -- (8.5,1.5) -- (8.5,2.5) -- (9.5,2.5);

  \node[font=\scriptsize\bfseries, ubGreen] at (0.5,0.5) {S};
  \node[font=\scriptsize\bfseries, ubRed] at (9.5,3.5) {G};
  \node[font=\scriptsize, ubMuted, align=left, anchor=west] at (10.4,3.2)
    {$\mathrm{g}$: biaya dari awal\\[2pt]
     $\mathrm{h}$: perkiraan ke tujuan\\[2pt]
     $f = g + h$};
\end{tikzpicture}

\end{center}

Gambar di atas memperlihatkan peta sepuluh kali enam petak. Petak gelap
adalah dinding, \textbf{S} titik awal, dan \textbf{G} tujuan. Petak berwarna
muda adalah jalur yang ditemukan; perhatikan bahwa jalur berbelok di sekitar
dinding, bukan menembusnya. Keterangan di sebelah kanan mengingatkan arti
$g$, $h$, dan $f$ yang baru saja Anda hitung.

\begin{ujibox}
Setelah Langkah 10, periksa lima hal berikut.
\begin{itemize}[leftmargin=1.4em, itemsep=0.2em]
  \item Penjaga berjalan memutar ketika pilar menghalangi garis lurus menuju
        pemain. Bila ia tetap menekan dinding, periksa pemanggilan
        \code{find_path} pada \code{_chase}.
  \item Garis kuning tidak pernah melewati petak dinding. Bila melewatinya,
        periksa penyaringan dinding pada \code{neighbours} dan
        \code{build_grid}.
  \item Garis kuning berakhir di petak tempat pemain berdiri. Bila berakhir
        di tempat lain, periksa urutan \code{_player_cell} dan pembalikan
        daftar pada \code{_reconstruct}.
  \item Penjaga sampai di ujung jalur dan berhenti tanpa bergetar. Bila
        bergetar, periksa \code{WAYPOINT_REACHED} dan pembatasan langkah
        pada \code{Guard.follow}.
  \item Ketika pemain berpindah petak, jalur dihitung ulang dan garis kuning
        berubah. Bila tidak berubah, periksa pembanding
        \code{last_player_cell} pada \code{_chase}.
\end{itemize}
\end{ujibox}

% ---------------------------------------------------------
\bagian{Kode Sesi Ini}

Potongan berikut adalah bagian pokok checkpoint sesi 8. Berkas lengkapnya,
termasuk kelas pemain dan level Rak yang tidak berubah, tersedia pada folder
\file{code/sesi-08}. Bandingkan kode Anda dengan daftar ini setelah program
berjalan, bukan sebelumnya.

\tampilkodebaris[Python]{code/sesi-08/gudang_artefak/settings.py}{44}{58}
  {Pengaturan petak, penjaga, dan pencarian jalur}

\tampilkodebaris[Python]{code/sesi-08/gudang_artefak/pathfinding.py}{27}{39}
  {Jarak Manhattan dan penentuan tetangga yang dapat dilalui}

\tampilkodebaris[Python]{code/sesi-08/gudang_artefak/pathfinding.py}{42}{83}
  {Pencarian A* dengan antrean prioritas dan rekonstruksi jalur}

\tampilkodebaris[Python]{code/sesi-08/gudang_artefak/game.py}{187}{202}
  {Menghitung ulang rute ketika petak pemain berubah}

\tampilkodebaris[Python]{code/sesi-08/gudang_artefak/game.py}{242}{264}
  {Menggambar dinding, pintu keluar, dan jalur hasil pencarian}

\checkpoint{08}{Program pada folder itu dapat dijalankan dengan
\code{python main.py --level 2}; penjaga menempuh jalan memutar di sekitar
pilar, dan jalur yang sedang dipakai tampil sebagai garis kuning beserta
titik pusat setiap petaknya.}

% ---------------------------------------------------------
\bagian{Kesalahan Umum}

\begin{galatbox}
\begin{itemize}[leftmargin=1.4em, itemsep=0.35em]
  \item \textbf{Jalur menembus dinding.}
        Gejala: garis kuning melintasi petak gelap. Penyebabnya, fungsi
        \code{neighbours} tidak menyaring dinding atau batas peta. Periksa
        syarat \code{0 <= nx < width}, \code{0 <= ny < height}, dan
        \code{(nx, ny) not in walls}.
  \item \textbf{\code{TypeError: '<' not supported between instances}.}
        Gejala: program berhenti saat antrean berisi dua petak dengan $f$
        sama. Penyebabnya, isi antrean hanya \code{(f, petak)} dan Python
        membandingkan objek yang tidak sejenis. Periksa susunan rangkap tiga
        \code{(f, g, petak)}.
  \item \textbf{Jalur yang ditemukan lebih panjang daripada seharusnya.}
        Gejala: penjaga memutar jauh padahal ada jalan dekat. Penyebabnya,
        isi antrean yang usang tidak dibuang, sehingga petak diperiksa
        dengan biaya lama. Periksa \code{closed} dan perbandingan
        \code{new_g < best_g.get(nxt, 10 ** 9)}.
  \item \textbf{Jalur terbalik.}
        Gejala: penjaga berjalan menjauhi pemain, lalu berbalik di ujung
        peta. Penyebabnya, daftar hasil rekonstruksi tidak dibalik. Periksa
        \code{path.reverse()} sebelum daftar dikembalikan.
  \item \textbf{Jalur kosong padahal tujuan terlihat.}
        Gejala: penjaga berhenti dan tidak ada garis kuning. Penyebabnya,
        petak awal atau tujuan ikut terhitung sebagai dinding, atau peta
        terkurung. Periksa keluaran \code{build_grid} dan bangun peta yang
        tujuan tidak dapat dicapai untuk memastikan program menanganinya.
  \item \textbf{Garis kuning muncul, tetapi penjaga diam.}
        Gejala: rute benar dan tergambar, badan penjaga tidak bergerak.
        Penyebabnya, \code{guard.path} tidak diisi atau \code{path_index}
        tidak dinolkan. Periksa kedua baris itu di dalam \code{_chase}.
\end{itemize}
\end{galatbox}

% ---------------------------------------------------------
\bagian{Lembar Kerja}

Isi tabel berikut pada catatan praktikum Anda. Lakukan pengukuran pada peta
yang sama agar angka dapat dibandingkan. Kolom terakhir diisi setelah Anda
mengubah satu nilai dan mengamati akibatnya.

{\small
\renewcommand{\arraystretch}{1.4}
\begin{tabularx}{\linewidth}{@{}C{0.8cm} L{6.0cm} Y@{}}
\toprule
\textbf{No} & \textbf{Hasil yang dicatat} & \textbf{Keterangan Anda} \\
\midrule
1 & Panjang jalur dari petak awal penjaga ke petak pemain, dalam petak & \\[0.55cm]
2 & Nilai $g$, $h$, dan $f$ untuk satu petak yang Anda pilih sendiri & \\[0.85cm]
3 & Banyak petak yang diperiksa sebelum tujuan ditemukan & \\[0.55cm]
4 & Apakah jalur melewati dinding, diperiksa satu per satu & \\[0.55cm]
5 & Nilai Kuis 2 (Materi 6--7) dan satu konsep FSM yang perlu diulang & \\[0.85cm]
6 & Satu kendala yang Anda hadapi dan cara menyelesaikannya & \\[0.9cm]
\bottomrule
\end{tabularx}
}

% ---------------------------------------------------------
\bagian{Checklist Sesi}

\begin{checklistbox}
\begin{checklist}
  \item Saya dapat menunjukkan penjaga yang menemukan jalan memutar di
        sekitar pilar.
  \item Saya dapat menjelaskan bagaimana peta berpetak menjadi graf, termasuk
        peran dinding dan biaya setiap sisi.
  \item Saya dapat menghitung $g$, $h$, dan $f$ untuk sebuah petak, serta
        menjelaskan arti setiap lambang.
  \item Saya dapat menjelaskan mengapa jarak Manhattan tidak melebihkan jarak
        sebenarnya pada gerak empat arah.
  \item Saya dapat menjelaskan cara \code{heapq} mengurutkan petak dan cara
        isi antrean yang usang dibuang.
  \item Saya dapat menjelaskan rekonstruksi jalur dari \code{came_from} dan
        menangani keadaan tujuan yang tidak dapat dicapai.
  \item Saya dapat membaca jalur yang digambar di layar dan memastikan jalur
        itu tidak menembus dinding.
\end{checklist}
\end{checklistbox}

% ---------------------------------------------------------
\bagian{Pengait ke Sesi Berikutnya}

Pencarian pada sesi ini dijalankan setiap kali petak pemain berubah. Cara itu
memadai untuk satu penjaga pada peta kecil, tetapi belum layak untuk mini
proyek. Pada sesi 9 pencarian dipanggil hanya ketika penjaga berada pada
keadaan mengejar, lalu diulang secara berkala setiap 0,3 detik. Kedua level
disusun berurutan dalam satu program, dan tiga skenario uji disiapkan untuk
menilai perilaku penjaga: lorong terbuka, lorong berbelok, dan halangan yang
memaksa jalur memutar.
